\section{Introduction}
\label{sec:introduction}

Robotic table tennis is a canonical benchmark for fast, contact-rich manipulation
because every rally compresses perception, prediction, planning, and execution
into a few hundred milliseconds. The agent must observe a fast-moving projectile,
forecast its bounce, position a 7-DOF arm, and impart a precise momentum
exchange through a small paddle face --- all while obeying joint limits and
respecting the geometric constraints imposed by the table and the net.

This term project addresses both phases of the ECE 275 robotic table tennis
benchmark. In the \emph{single-player} phase, the robot is served balls and
must return them past the net onto the opponent's half. In the
\emph{two-player} phase, two robots play complete matches against each other,
each subject to the same restrictions. The simulator is built on PyBullet and
ships a KUKA LBR iiwa 7 manipulator with a rigidly attached paddle; the
implementation lives in the \texttt{RobotPlayer} class inside
\texttt{one\_player\_game.py} and \texttt{two\_player\_game.py}.

We treat the problem as one of model-based planning rather than reactive
tracking. The agent uses the physics it can predict (gravity, restitution) to
\emph{lock} a swing plan as soon as the ball is heading its way, then executes
the plan open-loop in joint space. Inverse kinematics with explicit
null-space handling, a fixed strike plane, a time-scheduled
windup--follow-through sweep, and an empirically tuned set of geometric
parameters together produce a controller that wins every single-player game
in the benchmark and competes successfully in the two-player matches.

\paragraph{Report organisation.}
Section~\ref{sec:problem} states the problem formally. Section~\ref{sec:solution}
develops the algorithm. Section~\ref{sec:results} reports the simulator results.
Section~\ref{sec:discussion} discusses the insights and the connection to course
material, and Section~\ref{sec:conclusion} concludes.
